\section{Esquema de la fila CSV\_PEATONAL}
\label{anx:csv}
% Origen: filaResultadoPed() del modelo y CAMPOS de cargar_corridas.py (20 campos).

Cada corrida escribe una fila separada por punto y coma. El primer campo es el prefijo:
\texttt{CSV\_PEATONAL} (cargable), \texttt{CSV\_PEATONAL\_DEMO} (demostración, nunca se carga) o
\texttt{CSV\_PEATONAL\_INCOMPLETO} (corrida sin drenar, bloquea la carga).

\begin{longtable}{@{}r>{\raggedright\arraybackslash}p{5.2cm}>{\raggedright\arraybackslash}p{7.5cm}@{}}
\toprule
Campo & Nombre & Detalle \\
\midrule
\endhead
1 & Escenario & \EO{}, \EI{}, \ET{} o \EIT{} \\
2 & Semilla & Común a los cuatro escenarios de la réplica \\
3 & Pasajeros generados & Igual en los cuatro escenarios de la réplica \\
4 & Pasajeros desviados & Cero en todos los escenarios evaluados \\
5 & Procesados a las 09:30 & Validaciones hasta el corte \\
6 & Procesados con drenaje & Igual a los generados (conservación) \\
7 & Espera media & Segundos, desde la entrada a la cola hasta el inicio de la validación \\
8 & P90 de espera & Primaria \\
9 & Proporción con espera mayor a 30~s & Primaria, fracción entre 0 y 1 \\
10 & Lq & Largo medio de la cola en la franja \\
11 & Cola máxima & Personas \\
12 & Tiempo de drenaje final & Último egreso menos último ingreso \\
13 & Utilización media & Fracción, entre los molinetes habilitados \\
14 & Dispersión de utilización & Desvío entre molinetes \\
15 & Pasajeros de la cohorte pico & Ingresados de 08:15 a 08:45 \\
16 & Espera media de la cohorte pico & Segundos \\
17 & P90 de la cohorte pico & Primaria \\
18 & Día SBASE & Índice del día sorteado (0 a 74) \\
19 & P90 de tránsito del Roca a la cola & Segundos desde que baja del tren hasta que entra a la cola \\
20 & CV de validaciones entre molinetes & Para validar el reparto contra SBASE \\
\bottomrule
\end{longtable}

\section{Parámetros del modelo en la producción}
\label{anx:parametros}

\begin{table}[H]
\centering
\caption{Parámetros de \texttt{MainPeatonal} en \texttt{PeatonalCorridasApareadas}}
\label{tab:anx-parametros}
\small
\begin{tabularx}{\textwidth}{@{}>{\raggedright\arraybackslash}p{4.4cm} l X@{}}
\toprule
Parámetro & Valor & Significado \\
\midrule
\texttt{molinetesOperativosPed} & 20 o 22 & Molinetes habilitados \\
\texttt{servicioPedSeg} (min, máx) & 2,4 (1,8; 3,5) o 1,8 (1,5; 2,5) & Validación triangular: SUBE o sin contacto \\
\texttt{proporcionRocaPed} & 0,80 & Parte de las validaciones que llega en tren [supuesto] \\
\texttt{demoraAccesoRocaSeg} & 10~s & Del arribo a la apertura de puertas [supuesto] \\
\texttt{duracionDescargaSeg} & 30~s & Descenso por las 21 puertas [supuesto] \\
\texttt{diaSbasePed} & $-2$ & Día sorteado con la semilla entre los 75 días \\
\texttt{llegadasRealesPed} & verdadero & Trenes con la llegada real del día \\
\texttt{cargaTrenPoissonPed} & verdadero & Pasaje de cada tren con distribución de Poisson \\
\texttt{atrasoMaxRocaSeg} & 0 & Atraso adicional (solo para sensibilidad con el horario oficial) \\
Velocidad de marcha & U(1,1; 1,5)~m/s & Generador propio \\
Escaleras mecánicas & 0,685~m/s & 10,5~m con 4,6~m de desnivel a 0,30~m/s vertical \\
\bottomrule
\end{tabularx}
\end{table}

\section{Días hábiles y trenes del Roca}
\label{anx:roca}

Los 75 días del estudio son los días hábiles de marzo a junio de 2026 con datos de SBASE y servicio normal del
Roca. Se excluyeron cinco días con menos de 40 trenes en la franja, por servicio interrumpido: 23/03, 07/04,
13/04, 26/05 y 18/06. Las llegadas de cada tren a Plaza Constitución provienen de la respuesta de Trenes
Argentinos \parencite{sofse2026nota}; los perfiles por ventana, del dataset de SBASE. Los dos se embeben en el
modelo con \texttt{datos/embeber\_dias\_sbase.py}.
